\section{Equilibria, folds, and multistability for threshold networks}\label{s:KBmulti}

\subsection{Autonomous response networks undergo folds and cusps}
The derivative tests below are the standard scalar reduction of an
equilibrium equation with a simple zero eigenvalue
\cite{kuznetsov2004}. Theorem~\ref{t:Spe}(3) transfers these bifurcations from invariant
group reductions when the complementary directions are stable.
%Parameters are constant along trajectories.
On an invariant siphon face,
full stability additionally requires stability of the absent directions.

\beP[Equilibria admit saddle-node and cusp tests, and nondegenerate cusps persist]
[][sadCus]\label{p:cusprobust}
Let $f$ be $C^2$ near an equilibrium $(x_*,\theta_*)$ of
$\dot x=f(x;\theta)$, with $x\in\mathbb R^{n_S}$, $\theta\in\mathbb R^p$.
Assume $J=\partial_xf$ has a simple zero eigenvalue and all its other
eigenvalues have negative real parts.
A saddle-node bifurcation is also called a fold bifurcation.
Evaluate all derivatives below
at this equilibrium and choose
\[
 Jr=0,\quad \ell^\top J=0,\quad \ell^\top r=1,\qquad
 Q=I-r\ell^\top,\qquad \mathcal B=\partial_x^2f.
\]
\BEN \im
Varying $\theta_j$ gives a nondegenerate saddle-node if
\[
 a_j:=\ell^\top f_{\theta_j}\ne0,\qquad
 b_2:=\ell^\top\mathcal B[r,r]\ne0.
\]
\im
For the cusp test assume additionally that $f$ is $C^3$ and put
$\mathcal C=\partial_x^3f$.
If $b_2=0$, define uniquely
\[
 Jh_{20}=-\mathcal B[r,r],\quad\ell^\top h_{20}=0,
 \qquad Jh_j=-Qf_{\theta_j},\quad\ell^\top h_j=0.
\]
In this case, a nondegenerate cusp is unfolded by $(\theta_1,\theta_2)$ if
\[
 b_3:=\ell^\top\bigl(\mathcal C[r,r,r]
                  +3\mathcal B[r,h_{20}]\bigr)\ne0,
 \qquad
 \det\begin{pmatrix}a_1&a_2\\d_1&d_2\end{pmatrix}\ne0,
\]
where $d_j=\ell^\top\bigl((\partial_xf_{\theta_j})r+
\mathcal B[r,h_j]\bigr)$.
\im
Under the cusp conditions above, every sufficiently small $C^3$
perturbation of the vector-field family has a unique nearby cusp,
with shifted state and unfolding parameters $(\theta_1,\theta_2)$.
The cubic and unfolding conditions and the strictly stable complementary
spectrum persist. If the family depends smoothly on additional
perturbation parameters, the continued cusp depends smoothly on them
to the order supplied by the implicit-function theorem.
The continuation remains admissible when the original state and
unfolding parameters are interior to their admissible domains;
at a parameter boundary, admissibility must be checked separately.
\EEN
In all three assertions, parameters other than the two designated
unfolding parameters (or the single saddle-node parameter) are held fixed.
\eeP
\begin{proof}
The projection $Q$ has range $\ker\ell^\top$, which $J$ preserves,
and $J$ is invertible on that range. The implicit-function theorem
therefore solves
\[
 Qf(x_*+\xi r+w;\theta)=0,\qquad w\in\ker\ell^\top,
\]
uniquely as $w=w(\xi,\theta)$ near $(0,\theta_*)$. Put
\[
 g(\xi,\theta)=\ell^\top
 f(x_*+\xi r+w(\xi,\theta);\theta).
\]
Nearby equilibria correspond exactly to $g=0$. At the base point,
$w=0$, $w_\xi=0$, and differentiation of the projected equation gives
\[
 Jw_{\xi\xi}=-Q\mathcal B[r,r],\qquad
 Jw_{\theta_j}=-Qf_{\theta_j}.
\]
If $b_2=0$, then $Q\mathcal B[r,r]=\mathcal B[r,r]$, proving both
solvability and uniqueness of $h_{20}$; also $w_{\theta_j}=h_j$.
Since $\ell^\top J=0$, the scalar derivatives are
\[
 g_{\theta_j}=a_j,\qquad g_{\xi\xi}=b_2,
 \qquad g_{\xi\theta_j}=d_j,
 \qquad g_{\xi\xi\xi}=b_3\quad(b_2=0).
\]
These are the usual saddle-node and cusp tests for the equilibrium equation
\cite{kuznetsov2004}. The stable complementary spectrum gives the
corresponding differential-equation bifurcations. In particular, $g$
is obtained by solving the complementary equilibrium equations; the
center-manifold differential equation has the same nondegeneracy tests.

For persistence, apply the implicit-function theorem to
$\mathcal H=(g,g_\xi,g_{\xi\xi})$ in the variables
$(\theta_1,\theta_2,\xi)$. At the cusp,
\[
 D\mathcal H=
 \begin{pmatrix}a_1&a_2&0\\d_1&d_2&0\\ *&*&b_3\end{pmatrix},
 \qquad
 \det D\mathcal H=b_3
 \det\begin{pmatrix}a_1&a_2\\d_1&d_2\end{pmatrix}\ne0.
\]
The same complementary-equation reduction is available for nearby
families. Small $C^3$ perturbations give small $C^1$ perturbations
of $\mathcal H$, so its locally unique zero persists. Continuity
preserves the strict nondegeneracy and spectral conditions, and
interior admissibility is open.

For linear $P_i$ in \eqref{polR}, with $\varphi_i=\Phi_i'$,
the required state derivatives are
\[
 \mathcal B_i[r,r]=b_i\beta^2\varphi_i'(h_i+\beta(Wx)_i)(Wr)_i^2,
 \qquad
 \mathcal C_i[r,r,r]=b_i\beta^3\varphi_i''(h_i+\beta(Wx)_i)(Wr)_i^3.
\]
The correction terms use the same bilinear form:
\[
 \mathcal B_i[r,z]
 =b_i\beta^2\varphi_i'(h_i+\beta(Wx)_i)(Wr)_i(Wz)_i.
\]
\end{proof}

% p:KBinterpolate is retained as an alias of t:KBsharp.
% Elsewhere cite the interpolation construction in the proof of
% Theorem~\ref{t:KBsharp}, rather than a separate proposition.
